\documentclass[12pt]{article}
\usepackage[utf8]{inputenc}
\usepackage[T1]{fontenc}
\usepackage[margin=1in]{geometry}
\usepackage{setspace}
\usepackage{xcolor}
\usepackage[colorlinks=true, urlcolor=blue!60!black, linkcolor=blue!60!black]{hyperref}
\setlength{\parindent}{2em}
\setlength{\parskip}{1.25em}
\renewcommand{\baselinestretch}{1.0}
\begin{document}
\begin{center}
{\LARGE \textit{Dred Scott v. Sandford}, 60 U.S. 393 (1857)\par}
\vspace{0.8em}
{\large Case Brief --- Four-Page Edition\par}
\end{center}
\vspace{1em}

\section{Memory Jogger}

The Supreme Court held that people of African descent could not be citizens of the United States under the Constitution and struck down the Missouri Compromise of 1820, accelerating the sectional crisis that led to the Civil War before being repudiated by the Thirteenth and Fourteenth Amendments (\href{https://supreme.justia.com/cases/federal/us/60/393/}{Justia}).

\section{Detailed Case Facts}

Dred Scott, born enslaved in Virginia around 1799, was sold in St.\ Louis to Army surgeon Dr.\ John Emerson, who took him in 1834 to Rock Island in free Illinois and then to Fort Snelling in the Wisconsin Territory, where Congress had ``forever prohibited'' slavery in the Missouri Compromise. At Fort Snelling Scott married Harriet Robinson; their daughter Eliza was born north of Missouri. After Emerson died in 1843, his widow Irene continued to hire out the Scotts and refused Scott's 1846 offer to buy his family's freedom. The case unfolded amid the Kansas--Nebraska conflict, and President-elect Buchanan improperly pressed Justice Grier to join the Southern majority, promising in his inaugural that the question would be ``speedily and finally settled.''

\section{Procedural History}

In 1846 the Scotts sued for freedom in St.\ Louis under Missouri's ``once free, always free'' precedents and won a jury verdict in 1850, but the Missouri Supreme Court reversed 2--1 in \textit{Scott v. Emerson} (1852). Scott then sued John F. A. Sanford, a New York citizen, in federal court on diversity. Sanford's plea in abatement argued that a ``negro of African descent'' could not be a citizen; Judge Wells overruled it, but the jury found for Sanford on the merits. After argument and reargument, the Supreme Court on March 6, 1857, reversed and remanded with directions to dismiss for want of jurisdiction.

\section{Judicial Votes}

\textbf{7--2} against Scott; all nine Justices wrote separately.
\begin{itemize}
\item \textbf{Opinion of the Court:} Chief Justice Taney, joined in full by Justice Wayne.
\item \textbf{Concurrences:} Wayne; Nelson (Missouri law only); Grier; Daniel; Campbell; Catron.
\item \textbf{Dissents:} Justices McLean and Curtis.
\end{itemize}

\section{Holding}

(1) Persons of African descent whose ancestors were imported and sold as slaves, whether enslaved or free, are not and cannot become ``citizens'' under the Constitution and cannot invoke diversity jurisdiction. (2) The Missouri Compromise's ban on slavery north of $36^{\circ}30'$ is void (six Justices), because Congress may not exclude slave property from the territories and doing so violates Fifth Amendment due process. (3) Under \textit{Strader v. Graham}, Scott's Illinois residence did not free him once he returned to Missouri. Critics from Curtis to Lincoln called everything after the jurisdictional ruling \textit{obiter dictum}; notably, only three Justices endorsed the broad citizenship holding, while six invalidated the Compromise.

\section{Analysis of Opinions}

\textbf{Taney (Court).} Separating state from federal citizenship, Taney read 1776 and 1787 opinion, discriminatory statutes, and the Slave Trade and Fugitive Slave Clauses to exclude Black people from ``We the People,'' warning that citizenship would give them ``the full liberty of speech'' and the right ``to keep and carry arms.'' Reaching the merits, he confined the Territory Clause to lands held in 1787 and held, in an early form of substantive due process, that excluding slave property from a territory violated the Fifth Amendment---the first federal statute struck down since \textit{Marbury}.

\textbf{Concurrences.} Wayne gave ``unqualified assent,'' saying ``the peace and harmony of the country required the settlement.'' Nelson rested solely on Missouri's authority over status. Grier joined Nelson and Taney. Daniel's proslavery opinion called slaves pure property. Campbell read the Territory Clause as a power over public lands only. Catron found the plea waived and affirmed Congress's territorial power, but voided the Compromise under the Louisiana Purchase Treaty and state equality.

\textbf{McLean (dissenting).} Citizenship requires freedom and domicile, not the vote; slavery is ``a mere municipal regulation'' that ended once Scott entered free soil.

\textbf{Curtis (dissenting).} Free Black men voted in five of the thirteen States at ratification and so were among ``the people of the United States.'' Having found no jurisdiction, the Court's territorial ruling was extrajudicial; Congress had broad territorial power, and slavery rested on positive local law.

\section{Examples: Future Applications}

\begin{enumerate}
\item \textbf{Same-Side: Free Black Merchant in 1858.} A free Massachusetts merchant sues a New York shipper in federal court; under \textit{Dred Scott} the suit must be dismissed for lack of diversity.
\item \textbf{Same-Side: Territorial Ban on Slavery.} A Nebraska territorial legislature bans slavery under ``popular sovereignty''; under Taney's reasoning it is void, since Congress cannot delegate a power it lacks.
\item \textbf{Opposite-Side: Freedom Suit in Massachusetts.} An enslaved person brought to Massachusetts seeks habeas there; the holdings turned on Missouri's sovereignty after return, so the court could grant freedom.
\item \textbf{Opposite-Side: Birthright Citizenship Today.} A State denies citizenship to U.S.-born children of a racial minority; the Fourteenth Amendment's Citizenship Clause, enacted to overturn \textit{Dred Scott}, forbids it (\textit{Wong Kim Ark}).
\item \textbf{Fence-Sitter: Slaveholder in Transit Through New York.} New York frees slaves brought voluntarily into the State. Nelson reserved this \textit{Lemmon} problem: state sovereignty over status favors New York, but Taney's and Catron's property theories could, as Lincoln warned, nationalize slavery.
\end{enumerate}

\section{Critique}

\textbf{Universal condemnation.} \textit{Dred Scott} heads the anticanon---Chief Justice Hughes's ``greatest self-inflicted wound.'' Fehrenbacher and Finkelman document its white supremacy and historical distortion: Curtis proved free Black men voted in five States, and Taney's Territory Clause reading defied text and seventy years of practice. The attempt to settle slavery by judicial fiat wrecked the Court's credibility and hastened war.

\textbf{Originalism and substantive due process.} Critics of originalism cite Taney's insistence that the Constitution ``must be construed now as it was understood at the time of its adoption.'' Originalists answer that Curtis's dissent is the disciplined model and that the text called enslaved people ``persons.'' Justices Scalia and Thomas cite \textit{Dred Scott} as the original sin of substantive due process.

\section{Key Quotations}

\begin{quote}
``They had for more than a century before been regarded as beings of an inferior order, and altogether unfit to associate with the white race, either in social or political relations; and so far inferior, that they had no rights which the white man was bound to respect; and that the negro might justly and lawfully be reduced to slavery for his benefit.''
\end{quote}
\noindent --- Chief Justice Roger B. Taney, Opinion of the Court.

\begin{quote}
``And an act of Congress which deprives a citizen of the United States of his liberty or property merely because he came himself or brought his property into a particular Territory of the United States, and who had committed no offence against the laws, could hardly be dignified with the name of due process of law.''
\end{quote}
\noindent --- Chief Justice Roger B. Taney, Opinion of the Court.

\begin{quote}
``The truth is that citizenship, under the Constitution of the United States, is not dependent on the possession of any particular political or even of all civil rights, and any attempt so to define it must lead to error.''
\end{quote}
\noindent --- Justice Benjamin R. Curtis, Dissenting Opinion.

\end{document}
